\begin{frame}{Przyjęta metodyka prowadzi od wymagań do weryfikacji}
  \begin{enumerate}
    \item \highlight{Analiza} -- wymagania, dane wejściowe i ograniczenia.
    \item \highlight{Projekt} -- architektura rozwiązania i dobór narzędzi.
    \item \highlight{Realizacja} -- implementacja oraz integracja elementów.
    \item \highlight{Walidacja} -- scenariusze badań i kryteria oceny.
  \end{enumerate}
  \vspace{5mm}
  \imageplaceholder[0.22\textheight]{Miejsce na schemat procesu lub architekturę rozwiązania}
\end{frame}

\begin{frame}{Najważniejszy element rozwiązania}
  \begin{columns}[T,onlytextwidth]
    \begin{column}{0.57\textwidth}
      \imageplaceholder[0.43\textheight]{Miejsce na diagram, model, zrzut interfejsu albo fotografię prototypu}
    \end{column}
    \begin{column}{0.39\textwidth}
      \begin{itemize}
        \item Nazwij element i jego rolę.
        \item Wyjaśnij przepływ informacji.
        \item Wskaż kluczową decyzję projektową.
        \item Podaj najważniejsze ograniczenie.
      \end{itemize}
    \end{column}
  \end{columns}
\end{frame}

\begin{frame}[fragile]{Przykład kodu źródłowego}
\begin{lstlisting}[style=weail,language=Python]
def moving_average(samples, window):
    """Oblicz średnią kroczącą z próbek."""
    # Obsługa polskich znaków: ąćęłńóśźż
    return [
        sum(samples[i:i + window]) / window
        for i in range(len(samples) - window + 1)
    ]
\end{lstlisting}
  \vspace{2mm}
  \small Pokazuj wyłącznie fragment potrzebny do uzasadnienia decyzji technicznej.
\end{frame}
